\chapter{LITERATURE REVIEW}
\label{ch:literature}

\section{Introduction}

This chapter reviews the three bodies of work on which the thesis draws. These are
i) the materials-science literature on electrospinning and bead formation, which
defines what is measured and why, ii) the image-analysis literature on fibrous
microscopy, which describes the way in which the measurement is currently made, and
iii) the computer-vision literature on segmentation, from which the methods used
here are taken. The gap between them is stated at the end of the chapter.

\section{Theoretical Background}

\subsection{Electrospinning and the Formation of Bead Defects}

In electrospinning, a polymer solution is fed through a spinneret that is held at a
high potential relative to a collector. Above a threshold voltage the electrostatic
repulsion of the surface charge overcomes the surface tension of the pendant
droplet, which then deforms into a Taylor cone and emits a charged jet. The jet
undergoes a bending instability and stretches by orders of magnitude while the
solvent evaporates, leaving a solid fibre that is deposited as a randomly oriented
non-woven mat \cite{doshi1995electrospinning, huang2003review}.
A detailed account of the jet dynamics governing this process is given by Reneker
and Yarin \cite{reneker2008electrospinning}.

The morphology of the deposit is determined by a balance of forces, where surface
tension favours the break-up of a cylindrical thread into droplets while
viscoelastic stress and charge repulsion oppose it. At low viscosity, low charge
density or high surface tension the instability is not fully suppressed and the jet
collapses partially, producing ellipsoidal swellings along the fibres. The standard
account of this behaviour was established by Fong, Chun and Reneker
\cite{fong1999beaded}, who showed that higher viscosity and higher net charge
density result in fewer and more elongated beads, whereas higher surface tension
has the opposite effect. Bhardwaj and Kundu \cite{bhardwaj2010electrospinning}
review the way in which concentration, conductivity, voltage, flow rate and
tip-to-collector distance interact in practice.

Since the bead content responds so directly to these parameters, bead count and
bead size are used as the diagnostic output of a parameter sweep, and a recipe is
refined until the mat is bead-free or until the bead population is acceptably
small. Beads are also important physically, as they reduce the specific surface
area that makes these mats useful in filtration and catalysis, perturb the
pore-size distribution that governs permeability and act as stress concentrators
that lower the tensile strength.

\subsection{Scanning Electron Microscopy of Nanofibre Mats}

SEM is the standard characterisation tool for these materials. It resolves features
well below the diffraction limit of light microscopy and produces images in which
the contrast derives largely from surface topography, making it suitable for the
assessment of fibre morphology. Micrographs are typically acquired at several
magnifications, a low one in order to sample enough fibres for a representative
measurement and a higher one in order to resolve individual features. The
instrument software burns an information banner into the image, in which the scale
bar, the accelerating voltage, the working distance and the magnification are
recorded. The banner is convenient for a human reader and provides the means by
which pixel measurements are converted to physical units. It is not, however, part
of the specimen, and therefore it has to be excluded before any automated analysis.

\subsection{Semantic and Instance Segmentation}

Semantic segmentation assigns a class label to every pixel, whereas instance
segmentation additionally separates the pixels belonging to distinct objects of the
same class. This distinction is decisive in the present task. A semantic model that
labels every bead pixel correctly still produces a single connected component for a
cluster of touching beads, so that the derived count is too low while the derived
size distribution contains a spurious large object. Since the laboratory endpoint
is a count and a size distribution rather than an area fraction, an instance-level
output is required.

\section{Related Work}

\subsection{Classical Image Analysis for Fibrous Microscopy}

The established approach in the materials literature is a sequence of hand-designed
operations, namely contrast normalisation, thresholding, morphological cleaning and
a splitting step. Otsu's method \cite{otsu1979threshold} selects a global threshold
by maximising the between-class variance and remains the default in tools such as
Fiji \cite{schindelin2012fiji}. Touching objects are then separated by the watershed
transform \cite{vincent1991watersheds}, which is usually applied to the distance
transform of the binary mask.

For electrospun materials in particular, DiameterJ \cite{hotaling2015diameterj}
provides a validated open-source pipeline for fibre diameter and porosity from SEM
micrographs and is widely used. Its scope, however, is the fibre network, as it
characterises diameters and intersections rather than isolating discrete defects.

These pipelines have real advantages, since they are transparent, they require no
training data and their failure modes are interpretable. They depend, however, on
the foreground being separable by intensity, and in micrographs of a dense fibre mat
this assumption does not hold, as bright fibres cross the field at every angle and
occupy the same grey levels as the bead surfaces, so that a global threshold is not
able to isolate the beads. Watershed splitting has in addition a well-known tendency
to over-segment on noisy distance transforms, and this is the failure that corrupts
a count most directly.

\subsection{Encoder--Decoder Networks for Semantic Segmentation}

U-Net \cite{ronneberger2015unet} introduced the encoder--decoder architecture with
skip connections that has dominated biomedical segmentation since then. The
contracting path captures context while the expanding path recovers spatial
resolution, and the skip connections restore the high-frequency detail that is lost
to pooling. The architecture was designed explicitly for the regime of interest
here, having been demonstrated on training sets of a few dozen annotated images
together with heavy augmentation. Falk et al.\ \cite{falk2019unet} later packaged
it for cell counting and morphometry, thus confirming that it transfers to
quantification tasks.

U-Net nevertheless remains a semantic model. The standard remedy for instance
separation, which consists in predicting object boundaries as an auxiliary class and
using them to split the foreground, is effective but indirect, as its accuracy
depends on how reliably thin boundaries can be predicted.

\subsection{Instance Segmentation}

Faster R-CNN \cite{ren2015faster} established the two-stage detection paradigm, in
which a region proposal network scores candidate boxes and a second stage classifies
and refines them. Mask R-CNN \cite{he2017mask} extends this design with a parallel
branch that predicts a binary mask inside each proposal, and replaces the coarse RoI
pooling with RoIAlign in order to preserve spatial correspondence. Since each object
is handled inside its own region, touching objects are separated by construction
rather than by post-processing, and it is this property that makes the architecture
a natural choice for counting tasks. Feature pyramid networks \cite{lin2017feature}
supply the multi-scale representation that allows a single model to handle objects
of very different apparent size, which is directly relevant to a dataset spanning a
sixfold magnification range.

Several methods developed for microscopy exploit the shape regularity of the
targets. StarDist \cite{schmidt2018stardist} represents each object as a star-convex
polygon, which is well matched to roundish and convex objects and degrades
gracefully when they touch. Cellpose \cite{stringer2021cellpose} learns spatial
gradient flows towards object centres and follows them in order to recover the
instances, and it generalises across cell types without retraining. Bead defects are
approximately ellipsoidal and convex, so that these shape priors are a plausible fit
and are noted here as alternatives.

\subsection{One-Stage Detection}
\label{sec:onestage}

The two-stage design described above spends its first stage in proposing regions and
its second in refining them. One-stage detectors remove the proposal step and predict
the class and the box directly on a dense grid of locations, and it is this that
makes them fast. The formulation was introduced by YOLO \cite{redmon2016yolo}, whose
principal early weakness was its accuracy on the objects falling between grid cells,
that is small and clustered ones. Lin et al.\ \cite{lin2020focal} identified the
cause as the class imbalance of dense prediction, where the overwhelming majority of
the candidate locations are background, and addressed it with a loss that
down-weights easy negatives. The modern YOLO releases \cite{jocher2023ultralytics}
keep the one-stage structure but drop the anchor boxes, predicting object extent
directly and adding an optional branch that emits a mask per detection, which brings
the family into the same comparison as the region-based instance models.

The reason for testing such a model here is not speed, which this application does
not require, but the different inductive bias. A region-based model reasons about one
object at a time inside its own proposal, whereas a one-stage model reasons
about the whole frame at once and has to resolve neighbouring objects through
non-maximum suppression rather than by construction. For a dataset in which the
objects touch this is a substantive difference and not an implementation detail, and
Chapter~\ref{ch:results} measures in which direction it operates.

Two cautions from the literature carry directly into the way in which the result
should be read. First, the weakness of this family on small objects is documented
rather than incidental, and Section~\ref{sec:small_objects} quantifies it against
the other architectures included in this comparison. Second, a comparison against a
region-based model varies the backbone, the label assignment, the loss and the
augmentation schedule at the same time, so that it answers whether a different
family performs better on this data but not which component is responsible.

\subsection{Detection of Small Objects}
\label{sec:small_objects}

The objects counted in this work are small in the precise sense used by the COCO
benchmark, where an object is defined as \emph{small} when its area falls below
$32 \times 32 = 1024$ pixels. Under this convention 82.8\% of the particles
annotated here are small, rising to 97\% in the $500\times$ micrographs, while only
1.2\% are large (Section~\ref{sec:calibration}). This is not an incidental property
of the dataset, since it determines which results in the literature form a fair
point of comparison. Detector accuracy degrades sharply and systematically with
object size, a decline that is documented across architectures and datasets in the
survey of Liu et al.\ \cite{liu2021survey} and is attributed there to the small
number of pixels contributed by each object, to the loss of that signal through
successive downsampling, and to the coarse alignment between small objects and the
anchors or feature-map cells that are meant to capture them.

The clearest evidence for this degradation is obtained by holding the model and the
dataset fixed and splitting the score by object size alone. Table~\ref{tab:coco_size}
summarises this for five architectures evaluated on COCO, where every entry is taken
from the paper that introduced the method. The same pattern is observed for all of
them, as the small-object score lies between a seventh and a half of the large-object
score for the same model on the same images. The ratio also separates the
architectures in a way that the pooled AP does not. The single-stage detectors of
that generation lose the most, with YOLOv2 recovering only a seventh of its
large-object accuracy on small objects, whereas the two-stage detectors built on a
feature pyramid lose the least. This constitutes the empirical basis for preferring
a pyramid-based two-stage model in the present work.

\begin{table}[htbp]
    \centering
    \caption{Average precision on COCO decomposed by object size. Every row is the
    published result of the method's own paper, so that the comparison holds the
    dataset and the evaluation protocol fixed and varies only the architecture. The
    final column reports the ratio $\mathrm{AP}_S / \mathrm{AP}_L$, that is the
    fraction of a model's large-object accuracy that survives on small objects.
    Mask R-CNN is scored on masks and the others on boxes.}
    \label{tab:coco_size}
    \begin{tabular}{llrrrrr}
        \toprule
        Method & Backbone & AP & AP$_{50}$ & AP$_S$ & AP$_L$ & $\mathrm{AP}_S/\mathrm{AP}_L$ \\
        \midrule
        YOLOv2 \cite{redmon2016yolo, lin2020focal} & DarkNet-19 & 21.6 & 44.0 & 5.0  & 35.5 & 0.14 \\
        SSD513 \cite{lin2020focal}        & ResNet-101      & 31.2 & 50.4 & 10.2 & 49.8 & 0.20 \\
        Mask R-CNN \cite{he2017mask}      & ResNet-101-FPN  & 35.7 & 58.0 & 15.5 & 52.4 & 0.30 \\
        Faster R-CNN \cite{lin2020focal}  & ResNet-101-FPN  & 36.2 & 59.1 & 18.2 & 48.2 & 0.38 \\
        RetinaNet \cite{lin2020focal}     & ResNet-101-FPN  & 39.1 & 59.1 & 21.8 & 50.2 & 0.43 \\
        \bottomrule
    \end{tabular}
\end{table}

A second body of work studies domains in which nearly every object is small, and it
is these absolute numbers that set a realistic expectation for the present task.
Table~\ref{tab:small_object_lit} collects the studies whose stated aim, as given in
the abstract, is the detection or segmentation of small objects, together with the
figure reported in this thesis, so that the two can be read on the same scale.

The first row provides the plainest statement of the difficulty. Cheng et al.\
\cite{cheng2023soda} quote in their abstract a single detector on a single benchmark
scored by object size, where DyHead reaches 28.3 AP on small objects on COCO
test-dev against 50.3 on medium and 57.5 on large ones. Small objects are therefore
not a variant of the same problem, as they are roughly half as solvable with the
same model on the same images. The mechanism is geometric and is quantified in the
same paper, where a six-pixel displacement of a predicted box costs a small object
almost all of its overlap and reduces the IoU from 100\% to 32.5\%, while the same
displacement leaves a medium or a large object at 56.6\% or 71.8\% respectively. A
localisation error that is negligible for a large object is thus fatal for a small
one.

The dedicated benchmarks confirm this in absolute terms. On SODA-D, whose objects
average 20.3 pixels, a standard Faster R-CNN reaches 28.9 AP \cite{cheng2023soda}.
The aerial benchmark AI-TOD-v2 \cite{xu2022tiny} is more severe still, as its
objects average 12.7 pixels, which is comparable to the 14.3-pixel median of the
$500\times$ micrographs used here, and a standard Faster R-CNN reaches only 12.8 AP
on it while scoring 0.0 on the very-tiny band of 2--8 pixels. Even the method
proposed in that paper reaches 24.7 AP. TinyPerson \cite{yu2020tinyperson} defines
its task in the same terms, targeting persons below 20 pixels in large-scale images
and identifying the scale mismatch between the pre-training and the target data as
the cause of the loss. This is the same argument that the magnified-input experiment
of Section~\ref{sec:upsample_results} tests directly.

Microscopy studies report far higher scores, although on a materially easier
problem, since both \cite{monchot2021titanium} and \cite{bals2023nanoparticles}
image particles against a clean and uniform substrate rather than through an
occluding fibre network, and both of them report region-overlap measures such as
DICE or IoU rather than an average precision that also has to assign every instance
correctly.

\begin{table}[htbp]
    \centering
    \caption{Published results of studies that state small-object detection or
    segmentation as their aim, together with the result of this thesis on the same
    scale. Average precision is reported on the 0--100 scale used by those papers
    rather than on the 0--1 scale used elsewhere in this work. The metrics are not
    mutually comparable, since AP penalises every false positive and every missed
    instance whereas DICE and IoU measure region overlap once the objects have been
    found, and therefore the column reports the headline figure of each study under
    its own name.}
    \label{tab:small_object_lit}
    \begin{tabular}{p{3.0cm}p{2.9cm}p{2.4cm}p{3.2cm}}
        \toprule
        Study & Domain and object size & Method & Reported result \\
        \midrule
        Cheng et al. \cite{cheng2023soda} &
        COCO test-dev; by size band &
        DyHead &
        AP 28.3 small, 50.3 medium, 57.5 large \\
        \addlinespace
        Cheng et al. \cite{cheng2023soda} &
        Driving (SODA-D); mean 20.3 px &
        Faster R-CNN &
        AP 28.9, AP$_{50}$ 59.7 \\
        \addlinespace
        Xu et al. \cite{xu2022tiny} &
        Aerial imagery; mean 12.7 px &
        Faster R-CNN &
        AP 12.8, AP$_{50}$ 29.9, AP 0.0 on 2--8 px \\
        \addlinespace
        Xu et al. \cite{xu2022tiny} &
        Aerial imagery; mean 12.7 px &
        DetectoRS + NWD-RKA &
        AP 24.7, AP$_{50}$ 57.4 \\
        \addlinespace
        Monchot et al. \cite{monchot2021titanium} &
        TiO$_2$ agglomerates, SEM &
        Mask R-CNN, transfer learning &
        DICE 0.95 \\
        \addlinespace
        Bals and Epple \cite{bals2023nanoparticles} &
        Nanoparticles, SEM and STEM &
        UNet++ &
        IoU 0.92, $F_1$ 0.96 \\
        \addlinespace
        This work (Ch.~\ref{ch:results}) &
        Beads in SEM; median 14.3 px at $500\times$ &
        Mask R-CNN &
        AP 12.0, AP$_{50}$ 29.8 on masks \\
        \bottomrule
    \end{tabular}
\end{table}

The last row is the reason for placing the table here rather than in the results
chapter. The $\mathrm{AP}_{50}$ of 29.8 obtained in this work on masks lies within a
tenth of a point of what a standard Faster R-CNN achieves with boxes on AI-TOD-v2, a
benchmark built specifically for tiny-object detection and whose mean object size is
within two pixels of the median obtained here. This coincidence is not presented as
a claim of equivalence, since the datasets, the metrics and the task all differ, but
it does locate the result. A figure that appears to be a failure when compared with
pooled COCO numbers is an ordinary one when compared with the literature that shares
this object size.

Two lessons carry over into the methodology. The first is that a feature pyramid is
not optional at this object size, since it is the mechanism by which small objects
retain a representation at a resolution where they still occupy several cells of the
feature map. The second is that an average precision in the range reported in
Chapter~\ref{ch:results} has to be read against Tables~\ref{tab:coco_size} and
\ref{tab:small_object_lit} rather than against the pooled COCO figures that are
usually quoted, as the latter are dominated by objects an order of magnitude larger
than these.

\subsection{Prior Work at Epoka University}
\label{sec:epoka_prior}

The problem addressed here has a direct antecedent in the research of the host
institution, which for several years has concerned the counting and segmentation of
small, low-contrast objects in microscopy. The imaging modality is different, being
unstained brightfield rather than SEM, but the difficulties are the same ones that
recur throughout this thesis, namely objects that touch, backgrounds that are not
uniform, and too few annotated images to allow training from scratch.

Uka et al.\ \cite{uka2020faster} applied Faster R-CNN to the counting of cells in
low-contrast micrographs, motivated by the observation that manual counting is
repetitive and does not scale, and that unstained acquisition, which is chosen in
order to avoid the side effects of staining, makes the objects harder to separate
from the background. The combination of a counting endpoint with a detector that
isolates each object inside its own region is the same argument that is made in
Section~\ref{sec:small_objects} for the present work, and it is the reason why a
region-based detector rather than a purely semantic model is used here as the
primary method.

The same group has established that the preprocessing applied before a network sees
an image changes materially what that network is able to learn. Uka et al.\
\cite{uka2020preprocessing} report the effect of preprocessing choices on the
classification of microscopy images suffering from non-uniform illumination and low
contrast. The contrast measurements presented in Figure~\ref{fig:contrast} indicate
that the present dataset has the same character, with roughly eighteen grey levels
separating a particle from the surrounding mat, and for this reason preprocessing is
treated in
Chapter~\ref{ch:methodology} as an experimental variable rather than as a fixed
preliminary step.

A subsequent line of work covers the segmentation side. Uka et al.\
\cite{uka2021basis} set out the image-analysis basis for evaluating
cell--biomaterial interaction, while Polisi Duro et al.\ \cite{polisi2024detection}
apply U-Net architectures to detection, segmentation, counting and average-area
estimation, which are the same four endpoints reported in Chapter~\ref{ch:results}.
More recent work \cite{polisi2023pruning, polisi2025efficient} examines how far such
models can be compressed by pruning and quantisation, and how the choice of
optimiser and loss function affects segmentation quality, which bears on the
training configuration justified in Section~\ref{sec:training_config}. Alongside
this segmentation line, Polisi Duro et al.\ \cite{polisiduro2025optimizing} pair
three pretrained feature extractors with six classifiers for cell-type
classification, finding the deepest extractor and the most flexible classifier to be
the best combination.

% TODO(supervisor): a conference abstract by this group on segmenting nanobeads
% on nanofibres -- the same object in the same modality as this thesis, and so
% its closest antecedent -- was reported by a literature search and attributed to
% NanoBalkan 2024. It could not be confirmed: the conference's full abstract book
% was searched and contains no such abstract and none of those authors. Ask the
% supervisor whether the work exists and where, then cite it here. Until then
% nothing is claimed about it.

Table~\ref{tab:epoka_results} summarises the headline result of each of these
studies where it could be obtained.

\begin{table}[htbp]
    \centering
    \caption{Reported results of the prior work of the host institution on
    microscopy image analysis. Two rows are taken from the papers themselves, which
    are open access. Two report what the indexed abstract states and no more, since
    the full text is behind institutional access. One is marked \emph{not
    retrieved}, as neither its text nor an indexed abstract could be obtained. The
    counting accuracy in the first row is a secondary attribution taken from a paper
    that cites it. No figure reported here is approximated, and the incomplete rows
    are retained so that the gap remains visible rather than hidden.}
    \label{tab:epoka_results}
    \begin{tabular}{p{3.1cm}p{3.2cm}p{2.4cm}p{2.9cm}}
        \toprule
        Study & Task and data & Method & Reported result \\
        \midrule
        Uka et al. \cite{uka2020faster} &
        Cell counting, unstained brightfield &
        Faster R-CNN &
        Counting accuracy 78\%, as reported by \cite{ferreira2024classification} \\
        \addlinespace
        Uka et al. \cite{uka2020preprocessing} &
        Classification of 20\,000 cell images; test sets of 300 to 8000 &
        CNN, varied preprocessing &
        No headline figure in the abstract \\
        \addlinespace
        Polisi Duro et al. \cite{polisi2024detection} &
        Detection, segmentation, counting, mean area &
        U-Net &
        \emph{not retrieved} \\
        \addlinespace
        Polisi et al. \cite{polisi2023pruning} &
        Cell segmentation from small input data &
        U-Net, augmented, pruned and quantised &
        No numeric result in the abstract \\
        \addlinespace
        Polisi et al. \cite{polisi2025efficient} &
        Cell segmentation; 1088 crops of $256 \times 256$ &
        U-Net, BCE--Jaccard, RMSprop &
        IoU 0.90--0.93; precision 0.973; under 1\% lost to quantisation \\
        \addlinespace
        Polisi Duro et al. \cite{polisiduro2025optimizing} &
        Cell-type classification, three lines &
        VGG16 features, neural network &
        Accuracy 0.895, $F_1$ 0.896, AUC 0.981 \\
        \addlinespace
        \bottomrule
    \end{tabular}
\end{table}

% TODO(supervisor): question L1. Three rows of Table~\ref{tab:epoka_results} are
% still incomplete because the full texts are behind IEEE or institutional
% access. Ask the supervisor for the PDFs -- he is an author on all of them.
% \cite{polisi2024detection} matters most, since it reports a counting endpoint
% and is the closest published comparison to Section~\ref{sec:counting}. The 78%
% for \cite{uka2020faster} is currently taken from a paper that cites it, not
% from the paper itself; confirm it against the original.

One number in that table requires care, since it is the one that a reader will reach
for first. The segmentation work of this group reports an intersection over union
between 0.90 and 0.93, against the $\mathrm{AP}_{50}$ of 0.298 reported in
Chapter~\ref{ch:results}. The two values are not on the same scale and cannot be
subtracted. An IoU of that kind is a \emph{pixel-level} measure over a foreground
that occupies a large fraction of the frame, as it asks how well the predicted
region overlaps the true region while remaining indifferent to whether two touching
objects have been separated. Average precision is an \emph{instance-level} measure,
where every object has to be found, delineated and matched one-to-one, and where
every spurious detection is counted against the score. A model can therefore
maintain an IoU near 0.9 while merging half of the objects that it finds, and this
is exactly the failure exhibited by the U-Net baseline of
Section~\ref{sec:comparison}.

% Both figures are computed by code/evaluate.py and stored in results/metrics.json
% under loso.<method>.pixel_agreement, pooled over the 11 held-out micrographs.
When measured on the present dataset in the metric used by that group, the same
U-Net architecture reaches a pixel-level DICE of 0.667 and an IoU of 0.501, whereas
Mask R-CNN reaches 0.625 and 0.454 respectively. The gap with respect to 0.90 is
therefore not a gap between this work and that one, but a gap between cells that are
large, dense and imaged against a controlled background and beads that are small,
sparse and embedded in an occluding fibre network. The comparison is reported here
so that the two bodies of work can be read together, while
Section~\ref{sec:small_objects} provides the literature-wide version of the same
caution.

\subsection{Segmentation from Few Annotated Images}

Training on a handful of annotated micrographs is standard practice in microscopy,
and the countermeasures are well established, namely transfer learning from large
natural-image corpora such as ImageNet \cite{deng2009imagenet} or COCO
\cite{lin2014microsoft}, together with aggressive geometric and photometric
augmentation \cite{shorten2019survey}. Microscopy is more forgiving than most
domains in this respect, since flips and rotations are genuinely label-preserving,
as a micrograph, unlike a photograph of a natural scene, has no canonical
orientation.

Augmentation also introduces a hazard that is easy to overlook. If augmented copies
are generated before the data are partitioned, transformed versions of the same
source image can appear both in the training and in the test set, so that the test
set measures memorisation rather than generalisation. The same concern applies more
broadly to any hierarchical dataset in which several images derive from one physical
specimen, a point that is returned to in Chapter~\ref{ch:methodology}.

\subsection{Foundation Models for Segmentation}

The Segment Anything Model \cite{kirillov2023segment} is trained on a very large
mask corpus and produces class-agnostic masks from point or box prompts, with strong
zero-shot transfer. For a domain with only 11 annotated images this is attractive,
but SAM has no notion of what a bead is, and when prompted densely on a fibre mat it
segments fibres, inter-fibre voids and beads indiscriminately. Its realistic role in
the present context is therefore that of an annotation aid or a mask proposal
generator rather than a standalone detector.

\section{Evaluation of Instance Segmentation}

Detection quality is conventionally summarised by an average precision computed over
intersection-over-union thresholds, following the COCO protocol
\cite{lin2014microsoft}. For microscopy, metrics that penalise merges and splits
explicitly are more informative. The aggregated Jaccard index
\cite{kumar2017dataset} accumulates intersection and union over matched objects and
charges the unmatched ones to the denominator, while panoptic quality
\cite{kirillov2019panoptic} factors performance into a recognition and a
segmentation term. Since the laboratory endpoint is a count and a size distribution,
these segmentation metrics are reported here alongside the counting error and a
direct comparison of the predicted and the manual size distributions.

\section{Research Gap}

Three gaps emerge from the work reviewed above.

First, the image-analysis tooling available to the electrospinning community is
aimed at the fibre network, that is at diameter, orientation and porosity, rather
than at discrete defects. DiameterJ and comparable tools do not count beads, and
consequently the quantification of bead defects is still performed by hand.

Second, where deep segmentation has been applied to fibrous microscopy it has
usually been semantic and has reported pixel-level overlap. For a task whose output
is a count this is the wrong figure of merit, since it does not distinguish a model
that separates two touching beads from one that merges them, although only the
former counts correctly. Studies that report counting error against a manual ground
truth, and that compare instance-level with semantic-level modelling on the same
images, are scarce.

Third, and most consequential for the credibility of any result in this area,
evaluation protocols are frequently not matched to the structure of the data. Small
microscopy datasets are hierarchical, with many images per specimen and many
augmented copies per image, and a partition that ignores this structure allows the
same specimen, or the same source image, to appear on both sides of the split,
leaving the reported scores optimistic by an unknown margin.

This thesis addresses these gaps by treating bead quantification as an instance
segmentation problem, by evaluating it with the counting and sizing errors that
determine its usefulness in the laboratory alongside conventional segmentation
metrics, and by adopting a specimen-wise protocol whose effect relative to the naive
alternative is measured rather than assumed.

% Question 11 was answered on 2026-08-24: the supervisor asked for all of the
% group's Epoka publications to be reviewed and their results included, which is
% now Section~\ref{sec:epoka_prior} and Table~\ref{tab:epoka_results}. What
% remains is not a framing question but a retrieval one, recorded at the TODO
% above: five of those seven results are behind institutional access.

\section{Summary}

Bead defects are the standard diagnostic of a misconfigured electrospinning process
and are currently quantified by hand, since the existing automated tools target the
fibre network instead. Classical thresholding fails on these micrographs because
beads and fibres are not separable by intensity. Encoder--decoder networks segment
the bead class well but merge touching instances, thus corrupting precisely the
quantity of interest. Region-based instance segmentation separates instances by
construction and is therefore the approach adopted here. The following chapter sets
out the dataset, the calibration and the evaluation protocol in detail.
